\section{El programa en assembly}
\label{sec:programa}

El juego completo está en \code{sw/game.s}, reproducido en el
anexo~\ref{anx:game}. Ocupa 108 líneas de instrucciones, que el assembler expande a
114 palabras de máquina (456 de los 2048 bytes de la RAM).

\subsection{Estructura}

El programa se divide en un cuerpo principal con dos bucles anidados, partida y ronda,
y seis subrutinas (tabla~\ref{tab:subrutinas}).

\begin{enumerate}
  \item \code{\_start}: carga la base de periféricos en \code{s0} y pone los displays
        en \code{00}.
  \item \code{game}: espera que se suelten los botones, parpadea hasta que se apriete
        uno (\code{wait\_start}), toma ese instante como semilla del generador
        pseudoaleatorio y reinicia el contador de rondas (\code{a5}) y la suma de
        tiempos (\code{a4}).
  \item \code{round}: enciende los cuatro LEDs, espera \code{WAIT\_CYCLES}, espera
        botones sueltos, elige el objetivo (\code{pick\_target}), lo enciende y lee el
        contador.
  \item \code{wait\_press}: espera un botón y lee el contador otra vez. Si el botón es
        incorrecto salta a \code{wrong}, que escribe \code{0xEE} en los displays y
        vuelve a \code{round}. Si es correcto, convierte los ciclos a décimas, suma,
        muestra el tiempo e incrementa \code{a5}.
  \item Si \code{a5} $<$ \code{ROUNDS} vuelve a \code{round}. Si no, apaga los LEDs,
        deja ver el último tiempo durante \code{WAIT\_CYCLES}, muestra el promedio y
        vuelve a \code{game}, donde \code{wait\_start} lo mantiene en pantalla mientras
        los LEDs parpadean.
\end{enumerate}

\begin{table}[H]
\centering
\caption{Subrutinas de \code{game.s}. Todas son hoja: no llaman a otras.}
\label{tab:subrutinas}
\begin{tabularx}{\linewidth}{l l X}
\toprule
Subrutina & Entrada / salida & Qué hace \\
\midrule
\code{wait\_cycles} & \code{a0} = ciclos & espera \code{a0} ciclos leyendo el contador \\
\code{wait\_release} & -- & vuelve cuando los botones llevan \code{DEBOUNCE\_CYCLES} sin apretarse \\
\code{wait\_start} & sale \code{a0} = contador & parpadea 0101/1010 hasta que se apriete un botón \\
\code{pick\_target} & sale \code{a0} = LED one-hot & avanza el LFSR (\code{s1}) 32 pasos y elige un LED \\
\code{divu} & \code{a0}, \code{a1} $\to$ \code{a0} & división entera sin signo por restas sucesivas \\
\code{show\_number} & \code{a0} = número & muestra \code{a0} en decimal en los displays, máximo 99 \\
\bottomrule
\end{tabularx}
\end{table}

\subsection{Uso de registros}

RV32E tiene solo 16 registros, así que se asignaron de forma fija
(tabla~\ref{tab:registros}). Como ninguna subrutina llama a otra, el registro de
retorno \code{ra} nunca se sobrescribe antes de un \code{ret}, y el programa no
necesita stack: \code{sp}, \code{gp} y \code{tp} quedan sin usar.

\begin{table}[H]
\centering
\caption{Asignación de registros.}
\label{tab:registros}
\begin{tabular}{lll}
\toprule
Registro & Contenido & Vida \\
\midrule
\code{s0} & base de periféricos \code{0x80000000} & todo el programa \\
\code{s1} & estado del LFSR & una partida \\
\code{a2} & contador al encender el objetivo ($c_0$) & una ronda \\
\code{a3} & LED objetivo en one-hot & una ronda \\
\code{a4} & suma de décimas de las rondas correctas & una partida \\
\code{a5} & rondas correctas & una partida \\
\code{a0}, \code{a1} & argumentos y resultado de subrutinas & una llamada \\
\code{t0}--\code{t2} & temporales & una llamada \\
\bottomrule
\end{tabular}
\end{table}

\subsection{Aritmética sin multiplicación, división ni desplazamientos}
\label{sec:sin-mul}

\begin{itemize}
  \item \textbf{Multiplicar por 2} es sumar un número consigo mismo:
        \code{add x, x, x}. Para escribir las decenas en los 4 bits altos del display
        se multiplican por 16 con cuatro sumas seguidas.
  \item \textbf{Desplazar un bit a la izquierda} es lo mismo que multiplicar por 2. Así
        avanza el LFSR (sección~\ref{sec:lfsr}) y así se construye la máscara del LED
        objetivo: se parte de \code{1} y se duplica tantas veces como indique el índice.
  \item \textbf{Dividir} es contar cuántas veces cabe el divisor restándolo
        (\code{divu}). Es lento para cocientes grandes, pero los cocientes del juego son
        pequeños: 99 como máximo para un tiempo y 99 para el promedio.
  \item \textbf{Separar decenas y unidades} para los displays es otra división por
        restas: se resta 10 mientras se pueda, y lo que queda son las unidades.
  \item \textbf{Leer el bit 31} de un registro no requiere desplazarlo: el bit 31 es el
        bit de signo, así que vale 1 exactamente cuando el registro es negativo
        (\code{bge s1, zero, ...}).
\end{itemize}

\subsection{Generación pseudoaleatoria}
\label{sec:lfsr}

El LED objetivo se elige con un LFSR (\emph{linear feedback shift register}) de Galois
de 32 bits con el polinomio $x^{32} + x^{22} + x^{2} + x + 1$, que es de período
máximo: recorre los $2^{32}-1$ estados distintos de cero antes de repetirse
\cite{xapp052}. Un paso del LFSR es:

\[
  s \leftarrow
  \begin{cases}
    (s \ll 1) \oplus \texttt{0x00400007} & \text{si el bit 31 de } s \text{ es } 1, \\
    (s \ll 1) & \text{si no.}
  \end{cases}
\]

La constante \code{0x00400007} tiene en 1 los bits 22, 2, 1 y 0, que corresponden a
los términos del polinomio bajo $x^{32}$. El código~\ref{lst:lfsr} implementa el paso
sin desplazamientos, con las dos equivalencias de la sección anterior.

\begin{lstlisting}[language=rvasm, caption={Paso del LFSR y elección del LED en \code{pick\_target}.}, label={lst:lfsr}]
pt_step:
    add  t0, s1, s1                 # t0 = s1 << 1
    bge  s1, zero, pt_no_tap        # bit 31 era 0: sin realimentacion
    xor  t0, t0, t2                 # t2 = LFSR_TAPS
pt_no_tap:
    mv   s1, t0
    addi t1, t1, -1
    bnez t1, pt_step                # LFSR_STEPS = 32 veces

    andi t0, s1, 3                  # indice del LED: 0..3
\end{lstlisting}

\paragraph{Semilla.} Un LFSR es determinista: con la misma semilla produce siempre la
misma secuencia. Si la semilla fuera el contador leído en un punto fijo del programa
tras el reset, sería la misma en cada encendido y el juego repetiría la secuencia de
LEDs. Por eso el juego espera que el jugador apriete un botón para empezar
(\code{wait\_start}), y la semilla es el valor del contador en ese instante. Un humano
no puede repetir un instante con precisión de 40~ns, así que cada partida parte de
un estado distinto. Se le aplica un \code{ori} con 1 para que la semilla nunca sea
cero, el único estado del que un LFSR no sale.

\paragraph{Por qué 32 pasos por ronda.} No todos los bits de la semilla son
aleatorios. Los bits altos del contador cambian lento: el bit 30 vale 0 durante los
primeros 43 segundos tras el encendido. Los bits más bajos tampoco son libres, porque
el bucle que lee el botón lo hace cada una cantidad fija de ciclos. Si el LFSR avanzara
pocos pasos por ronda, los dos bits que eligen el LED vendrían de los bits poco
aleatorios de la semilla. Avanzando 32 pasos, una palabra completa, cada bit del
estado pasa a depender de los bits medios de la semilla, que sí dependen de cuándo
apretó el jugador.

Esto se comprobó con un modelo del LFSR en Python (\code{sim/lfsr\_modelo.py}), que
simula 20.000 partidas con semillas como las que produce el programa: apretones entre
0,5 y 5 segundos tras el encendido, leídos cada $P$ ciclos
(tabla~\ref{tab:lfsr}). Con 2 pasos por ronda, el primer LED es \emph{siempre} el
mismo. Con 32 pasos, cada LED sale un 25\,\% de las veces, los 16 pares posibles de LEDs
consecutivos aparecen con la misma frecuencia, y el número de secuencias distintas
coincide con el que se espera de un generador ideal: con $4^{10} \approx 10^6$
secuencias posibles, en 20.000 partidas se esperan unas 190 repeticiones por azar.

\begin{table}[H]
\centering
\caption{Modelo de 20.000 partidas. $P$: cada cuántos ciclos lee el programa los botones.}
\label{tab:lfsr}
\begin{tabular}{ccccc}
\toprule
Pasos por ronda & $P$ & Primer LED: 0 / 1 / 2 / 3 (\%) & Pares distintos & Secuencias distintas \\
\midrule
2 & 1 & 100 / 0 / 0 / 0 & 1 de 16 & 14\,133 \\
2 & 16 & 100 / 0 / 0 / 0 & 1 de 16 & 14\,218 \\
32 & 1 & 25,5 / 24,9 / 25,0 / 24,6 & 16 de 16 & 19\,809 \\
32 & 16 & 24,7 / 25,2 / 25,1 / 25,0 & 16 de 16 & 19\,796 \\
\bottomrule
\end{tabular}
\end{table}

Cada paso toma 12 o 13 ciclos, según si hay realimentación, así que avanzar 32 pasos
cuesta unos 400 ciclos, 16~$\mu$s. Ocurren antes de encender el objetivo y por lo tanto
no afectan la medición.

\subsection{Instrucciones utilizadas}

La tabla~\ref{tab:instrucciones} lista las instrucciones de máquina que genera el
assembler para \code{game.s}, después de expandir las pseudo-instrucciones. El programa
usa 17 instrucciones distintas de RV32I y ninguna fuera de RV32E.

\begin{table}[H]
\centering
\caption{Instrucciones de máquina de \code{game.hex} (114 en total).}
\label{tab:instrucciones}
\begin{tabularx}{\linewidth}{l c X}
\toprule
Instrucción & Veces & Uso principal \\
\midrule
\code{addi} & 32 & constantes (\code{li}), copias (\code{mv}), contadores \\
\code{jal} & 18 & llamadas (\code{call}) y saltos (\code{j}) \\
\code{lw} & 12 & leer el contador de ciclos y los botones \\
\code{add} & 9 & multiplicar por 2 (máscara, LFSR, decenas) y acumular \\
\code{lui} & 7 & parte alta de constantes grandes (\code{li}) \\
\code{sub} & 6 & intervalos de tiempo y divisiones por restas \\
\code{bltu} & 6 & comparaciones sin signo de tiempos y cocientes \\
\code{jalr} & 6 & retorno de subrutina (\code{ret}) \\
\code{sw} & 6 & escribir LEDs y displays \\
\code{beq}, \code{bne} & 3, 3 & \code{beqz}/\code{bnez} y comparación de botón \\
\code{bge}, \code{bgeu} & 1, 1 & bit 31 del LFSR; saturación a 99 (\code{bleu}) \\
\code{andi}, \code{ori} & 1, 1 & índice del LED; semilla distinta de cero \\
\code{xor}, \code{xori} & 1, 1 & realimentación del LFSR; alternar el parpadeo \\
\bottomrule
\end{tabularx}
\end{table}
